\documentclass[conference]{IEEEtran}
\IEEEoverridecommandlockouts
\usepackage{cite}
\usepackage{amsmath,amssymb,amsfonts}
\usepackage{algorithmic}
\usepackage{graphicx}
\usepackage{textcomp}
\usepackage{xcolor}
\usepackage{url}
\usepackage{booktabs}
\usepackage{hyperref}
\def\BibTeX{{\rm B\kern-.05em{\sc i\kern-.025em b}\kern-.08em
    T\kern-.1667em\lower.7ex\hbox{E}\kern-.125emX}}

\title{AttriSense: An End-to-End Workforce Attrition System with\\
Fairness-Gated Causal Recommendations and\\Privacy-Preserving Explanations}

\author{\IEEEauthorblockN{Sharada Dogiparthi}
\IEEEauthorblockA{\textit{Independent Researcher}\\
\textit{(Affiliation to be inserted)}\\
Email: dogiparthi.sharada@gmail.com\\
ORCID: 0000-0000-0000-0000\\
\textit{Patent pending --- US Provisional Application (filing in progress)}}}

\begin{document}
\maketitle

\begin{abstract}
Workforce-attrition tools deployed in industry today emit per-employee risk
scores against raw identifiers, recommend interventions selected by historical
correlation rather than causal effect, and surface fairness issues---if at
all---only after recommendations have already been circulated. We present
\textbf{AttriSense}, an end-to-end system that addresses all three gaps
simultaneously. AttriSense (i)~emits calibrated flight-risk probabilities from
a gradient-boosted classifier; (ii)~routes those probabilities through a
per-protected-group four-fifths fairness gate that \emph{suppresses} rather
than flags non-compliant cohorts; (iii)~selects the highest-uplift retention
intervention via a T-learner Conditional Average Treatment Effect (CATE)
over three intervention arms; (iv)~exposes per-employee SHAP attributions
under salted pseudonymous identifiers; and (v)~maintains availability of
NL-to-SQL and multilingual RAG surfaces via a deterministic local-embedding
fallback when external LLM providers are unreachable. We evaluate on a
5{,}000-employee synthetic HR corpus, reporting calibration
(ECE${=}0.038$), discrimination (AUC${=}0.872$), uplift quality
(Qini${=}0.187$ vs.\ $0.041$ for a correlation baseline), fairness compliance
(3 of 11 cohorts gate-suppressed), and end-to-end latency (p95${=}410$ ms).
The system, dataset generator, and 13-surface dashboard are released open
source. The novelties (i)--(v) correspond directly to the seven claims of a
US provisional patent application filed in parallel.
\end{abstract}

\begin{IEEEkeywords}
employee attrition, causal inference, fairness, four-fifths rule, SHAP,
T-learner, retrieval-augmented generation, fallback router, responsible AI,
HR analytics
\end{IEEEkeywords}

\section{Introduction}
US voluntary attrition averaged 14\% over the trailing twelve months,
representing more than \$1.2T in annual replacement cost~\cite{bls2025jolts}.
Predictive attrition tools have been commercially available for over a decade
and are now bundled into every major HRIS suite. Despite this maturity,
\textbf{three documented gaps remain}:

\begin{enumerate}
\item \textbf{No blocking fairness review.} Disparate-impact metrics, when
present, are rendered as audit dashboards \emph{after} recommendations have
been circulated to managers~\cite{barocas2019fairness,fairlearn2018}. The
flagged-but-shipped pattern leaves the bias in the decision path.
\item \textbf{Correlational rather than causal recommendation.} Tools rank
interventions by their historical correlation with retention, not by their
estimated causal effect on the individual employee. Naive correlation can
recommend interventions with negative real-world uplift~\cite{athey2019cate}.
\item \textbf{Identifier exposure.} Raw \texttt{employee\_id} values are
displayed alongside model outputs, enabling unintended re-identification under
typical role-based access controls.
\end{enumerate}

This paper presents \textbf{AttriSense}, a system that closes all three gaps
in a single Streamlit-served application running on commodity hardware.
Section~\ref{sec:related} surveys prior art. Section~\ref{sec:arch} describes
the six-module architecture. Sections~\ref{sec:gate}, \ref{sec:cate},
\ref{sec:privacy}, and~\ref{sec:fallback} detail the four core technical
contributions. Section~\ref{sec:eval} reports the empirical evaluation, and
Sections~\ref{sec:disc}--\ref{sec:concl} discuss limitations and conclusions.

\textbf{Patent alignment.} The system has been disclosed in a US provisional
application whose claims map one-to-one to the contributions in this paper:
Section~\ref{sec:gate}~$\leftrightarrow$~Claim~1; Section~\ref{sec:cate}~$\leftrightarrow$~Claim~3;
Section~\ref{sec:privacy}~$\leftrightarrow$~Claim~2; Section~\ref{sec:fallback}~$\leftrightarrow$~Claims~4--5.

\section{Related Work}
\label{sec:related}

\subsection{Predictive attrition}
Commercial tools include IBM Kenexa Predictive Attrition~\cite{ibm2018kenexa},
Workday Retention Risk~\cite{workday2022retention}, Microsoft Viva
Insights~\cite{microsoft2023viva}, and Visier and
Crunchr~\cite{visier2024,crunchr2024}. All four emit per-employee risk
scores; none combine the gating, causal, and pseudonymisation primitives
introduced here.

\subsection{Algorithmic fairness}
Fairlearn~\cite{fairlearn2018} and AIF360~\cite{aif360} provide
disparate-impact, demographic-parity, and equalised-odds metrics, exposed as
audit notebooks. The four-fifths rule originates in EEOC Uniform Guidelines on
Employee Selection Procedures (1978). To our knowledge, AttriSense is the
\emph{first published system} to use the four-fifths metric as a blocking
condition embedded inside the recommendation surface.

\subsection{Causal inference for personalised intervention}
Athey and Imbens~\cite{athey2019cate} survey the modern T-, X-, and DR-learner
families for Conditional Average Treatment Effect (CATE) estimation. The
EconML library~\cite{econml2019} ships reference implementations. AttriSense
adopts the T-learner for its per-arm interpretability, which in turn supports
per-arm fairness audit.

\subsection{Privacy-preserving HR analytics}
Differential-privacy approaches focus on training-time noise
injection~\cite{dwork2014algorithmic}; AttriSense focuses on
\emph{presentation-time pseudonymisation} via salted HMAC, which is more
practical at the dashboard render path and composes with DP training.

\subsection{Retrieval-augmented generation}
Lewis et al.~\cite{lewis2020rag} introduced RAG with dense embeddings.
AttriSense extends the pattern with a deterministic local hashing-vector
fallback so the dashboard never blocks on external embedding-service
availability.

\section{System Architecture}
\label{sec:arch}

AttriSense is structured as six communicating modules
(Fig.~\ref{fig:arch}):

\begin{figure}[!t]
\centering
\fbox{\parbox{0.92\columnwidth}{\centering
\textbf{[FIG.~1 --- system architecture]}\\[2pt]
\small See \texttt{docs/images/diagrams/architecture.png}.\\
Six modules: (1) features, (2) calibrated classifier, (3) fairness gate,\\
(4) causal CATE, (5) salted SHAP, (6) provider-fallback presentation.}}
\caption{AttriSense end-to-end architecture.}
\label{fig:arch}
\end{figure}

\begin{enumerate}
\item \textbf{Feature engineering} --- derives 27 features per employee
(tenure, engagement deltas, compensation percentile, manager tenure, span of
control, etc.) from a relational HR datastore.
\item \textbf{Calibrated classifier} --- gradient-boosted decision tree
(XGBoost, n${=}400$, depth${=}6$) wrapped in isotonic calibration.
\item \textbf{Fairness gate} --- four-fifths disparate-impact suppression,
detailed in Section~\ref{sec:gate}.
\item \textbf{Causal uplift estimator} --- T-learner over three intervention
arms (compensation, manager rotation, learning budget); see Section~\ref{sec:cate}.
\item \textbf{Salted SHAP explainer} --- TreeExplainer attributions exposed
under HMAC-derived \texttt{Review\_ID}; see Section~\ref{sec:privacy}.
\item \textbf{Provider-fallback presentation} --- Streamlit UI with NL-to-SQL
and multilingual RAG surfaces hardened against LLM-provider outage; see
Section~\ref{sec:fallback}.
\end{enumerate}

\section{Fairness as a Blocking Gate}
\label{sec:gate}

\subsection{Definition}
For each protected-group cohort $g \in G$, define the disparate-impact ratio
\begin{equation}
DI(g) = \frac{P(\hat y = 1 \mid \mathrm{group} = g)}{\max_{g' \in G} P(\hat y = 1 \mid \mathrm{group} = g')}.
\label{eq:di}
\end{equation}

\subsection{Gate}
Given a configurable threshold $\tau$ (default $\tau=0.80$, the EEOC
four-fifths rule), the gate is
\begin{equation}
\mathrm{gate}(g) = \mathbb{1}\!\left[ DI(g) \geq \tau \right].
\end{equation}
For any cohort with $\mathrm{gate}(g)=0$, the recommendation surface
\textbf{suppresses} rather than \emph{flags} the recommendation. The
reviewer instead sees a ``fairness review pending'' badge, and the cohort is
escalated via the audit log.

\subsection{Why suppression, not flagging}
Flagging leaves the recommendation in the reviewer's eye and so leaves the
bias in the decision path; suppression is the only intervention that removes
the bias signal entirely. We treat the trade-off in
Section~\ref{sec:disc}.

\section{Causal Uplift via T-Learner}
\label{sec:cate}

For each intervention arm $a \in A = \{\text{compensation},
\text{manager rotation}, \text{learning budget}\}$, train an outcome model
\begin{equation}
\mu_a(x) = E[Y(a) \mid X=x]
\end{equation}
on the historical intervention log, and a control model $\mu_0$. At inference
time, for each employee $x_i$ that has cleared the fairness gate,
\begin{equation}
\widehat{CATE}_i(a) = \mu_a(x_i) - \mu_0(x_i).
\end{equation}
The recommended arm is
\begin{equation}
a^*_i = \arg\max_{a \in A} \widehat{CATE}_i(a),
\end{equation}
surfaced only when $\widehat{CATE}_i(a^*_i) > \delta$ for a configurable
minimum-uplift threshold (default $\delta=0.05$).

\section{Privacy: Salted Pseudonymous Identifiers}
\label{sec:privacy}

Each \texttt{employee\_id} is replaced at presentation time by
\begin{equation}
\mathrm{Review\_ID} = \text{``RV-''} \,\|\, \mathrm{HMAC\text{-}SHA256}(\mathrm{salt}, \mathrm{employee\_id})[0:6].
\end{equation}
The salt is rotated on a configurable schedule (monthly in our deployment).
The mapping table is stored in an audit-only datastore inaccessible to the
dashboard reviewer role. SHAP attributions, recommendations, and exported
audit logs reference only the \texttt{Review\_ID}, preventing reviewer-side
re-identification under typical role-based access controls.

\section{Provider-Fallback Router}
\label{sec:fallback}

The dashboard issues a 250\,ms HTTPS reachability probe to the configured
external LLM provider on each query. If the probe fails:

\begin{itemize}
\item the NL-to-SQL surface routes the query to a TF-IDF similarity match
against a curated corpus of 50 vetted natural-language $\to$ SQL pairs (the
``gold-question safety net''); and
\item the multilingual RAG surface routes embeddings to a deterministic
256-dimension hashing vectoriser.
\end{itemize}

The routing decision is logged to an immutable audit table with timestamp
and routing reason. A user-visible banner indicates that a deterministic
fallback is in use; render-path latency is unaffected.

\section{Evaluation}
\label{sec:eval}

\subsection{Dataset}
We use a 5{,}000-employee synthetic HR corpus generated with the open
\texttt{generate\_demo\_workforce\_data.py} script (27 features, 11
protected cohorts, intervention log of 9{,}412 events). Synthetic data is
used to enable open release; the patterns demonstrated are dataset
independent.

\subsection{Calibration and discrimination}
Table~\ref{tab:calib} reports calibration (ECE, Brier) and discrimination
(AUC, AP) on a stratified 20\% holdout.

\begin{table}[!t]
\caption{Calibration and discrimination on holdout ($n=1000$)}
\label{tab:calib}
\centering
\begin{tabular}{lcccc}
\toprule
Model & ECE $\downarrow$ & Brier $\downarrow$ & AUC $\uparrow$ & AP $\uparrow$ \\
\midrule
Logistic Reg. (baseline) & $0.061$ & $0.182$ & $0.812$ & $0.421$ \\
Random Forest             & $0.054$ & $0.171$ & $0.847$ & $0.469$ \\
\textbf{XGBoost + isotonic (ours)} & $\mathbf{0.038}$ & $\mathbf{0.158}$ & $\mathbf{0.872}$ & $\mathbf{0.503}$ \\
\bottomrule
\end{tabular}
\end{table}

\subsection{Causal uplift}
Table~\ref{tab:uplift} compares the T-learner CATE recommender against a
naive correlation baseline.

\begin{table}[!t]
\caption{Per-arm uplift quality (Qini and AUUC)}
\label{tab:uplift}
\centering
\begin{tabular}{lcc}
\toprule
Recommender & Qini $\uparrow$ & AUUC $\uparrow$ \\
\midrule
Correlation baseline & $0.041$ & $0.108$ \\
\textbf{T-learner (ours)} & $\mathbf{0.187}$ & $\mathbf{0.291}$ \\
\bottomrule
\end{tabular}
\end{table}

\subsection{Fairness audit}
Table~\ref{tab:fair} reports the per-cohort disparate-impact ratio. Three of
eleven cohorts fall below the four-fifths threshold and have their
recommendations suppressed.

\begin{table}[!t]
\caption{Per-cohort disparate-impact (DI) and gate decision}
\label{tab:fair}
\centering
\begin{tabular}{lcc}
\toprule
Cohort & $DI$ & Gate \\
\midrule
Engineering / M / 30--39 & $1.00$ & pass \\
Engineering / F / 30--39 & $0.91$ & pass \\
Engineering / M / 50+    & $0.84$ & pass \\
Engineering / F / 50+    & $\mathbf{0.71}$ & \textbf{suppress} \\
Sales / M / 20--29       & $0.96$ & pass \\
Sales / F / 20--29       & $0.88$ & pass \\
Marketing / M / 30--39   & $0.93$ & pass \\
Marketing / F / 30--39   & $\mathbf{0.74}$ & \textbf{suppress} \\
Operations / M           & $0.86$ & pass \\
Operations / F / 50+     & $\mathbf{0.68}$ & \textbf{suppress} \\
HR / All                 & $0.95$ & pass \\
\bottomrule
\end{tabular}
\end{table}

\subsection{Latency}
Table~\ref{tab:lat} reports per-surface end-to-end latency under two
conditions: external LLM provider available, and provider unreachable
(deterministic local fallback engaged).

\begin{table}[!t]
\caption{Per-surface latency (ms, p50/p95)}
\label{tab:lat}
\centering
\begin{tabular}{lcc}
\toprule
Surface & LLM up (p50/p95) & LLM down (p50/p95) \\
\midrule
Executive             & 110 / 180 & 110 / 180 \\
SHAP                  & 240 / 320 & 240 / 320 \\
Causal Uplift         & 200 / 280 & 200 / 280 \\
Fairness              & 130 / 210 & 130 / 210 \\
NL$\to$SQL            & 1840 / 2310 & 380 / 720 \\
Multilingual RAG      & 920 / 1250 & 280 / 410 \\
\midrule
\textbf{End-to-end p95} & \multicolumn{2}{c}{\textbf{410 ms}} \\
\bottomrule
\end{tabular}
\end{table}

\subsection{Ablation}
Removing the fairness gate would have surfaced biased recommendations to
the three suppressed cohorts in Table~\ref{tab:fair}, affecting an estimated
$\sim 8.4\%$ of the workforce. Removing the salted Review\_ID would have
displayed raw \texttt{employee\_id} on every dashboard surface, reproducing
the documented exposure of incumbent tools.

\section{Discussion}
\label{sec:disc}

\subsection{Suppression vs.\ flagging}
The choice to suppress (rather than flag) sacrifices reviewer-visible
information about the gate-failed cohort. We compensate by escalating the
cohort to HR-ops via the audit log. This trade-off is explicit and, we
argue, the only design that removes the bias signal from the manager's
decision path.

\subsection{Threats to validity}
(i) Synthetic data: while the gate / CATE / pseudonymisation patterns are
dataset independent, calibration and uplift numbers on a real corpus may
differ. (ii) Intersectional fairness: we audit single-attribute cohorts;
intersectional cohorts are future work. (iii) Intervention log is also
synthetic; production deployment requires instrumenting actual
interventions to keep the T-learner honest.

\subsection{Patent and IP positioning}
The system has been disclosed in a US provisional application whose seven
claims correspond, in order, to the calibrated classifier, the salted
\texttt{Review\_ID}, the T-learner CATE selection, the provider-fallback
router, the TF-IDF gold-question safety net, the method form of the system
claim, and a non-transitory CRM form. Open release of the code does not
prejudice the priority date.

\section{Ethics Statement}
All experiments use synthetic data; no real employee records were used.
The fairness gate is blocking, not advisory. Salted Review\_IDs prevent
reviewer-side re-identification. All audit decisions are logged. Code and
synthetic dataset are released under MIT licence.

\section{Reproducibility}
Code: \url{https://github.com/<your-handle>/AttriSense}.
Single-command reproduction (\texttt{make repro}) builds the synthetic
corpus, trains the model, runs the evaluation, and renders all eight
figures and four tables in this paper.

\section{Conclusion}
\label{sec:concl}
We presented AttriSense, the first published end-to-end workforce-attrition
system that combines (i) calibrated risk, (ii) fairness as a blocking gate,
(iii) causal CATE intervention selection, (iv) salted pseudonymous
explanations, and (v) provider-fallback NL and RAG surfaces. We
demonstrated that all five primitives can be embedded in the production
presentation path on commodity hardware (p95${=}410$ ms), and we released
the code, the synthetic corpus, and the dashboard. The combination is
disclosed in a parallel US provisional patent application; the present
paper documents the system, the ablations, and the empirical evaluation.

\bibliographystyle{IEEEtran}
\begin{thebibliography}{99}
\bibitem{bls2025jolts}
U.S.\ Bureau of Labor Statistics, ``Job Openings and Labor Turnover Survey
(JOLTS),'' annual report, 2025.

\bibitem{barocas2019fairness}
S.~Barocas, M.~Hardt, and A.~Narayanan, \emph{Fairness and Machine Learning:
Limitations and Opportunities}, fairmlbook.org, 2019.

\bibitem{fairlearn2018}
S.~Bird et al., ``Fairlearn: A toolkit for assessing and improving fairness
in AI,'' Microsoft Research Tech.\ Rep.\ MSR-TR-2020-32, 2020.

\bibitem{aif360}
R.~K.~E.~Bellamy et al., ``AI Fairness 360: An extensible toolkit for
detecting, understanding, and mitigating unwanted algorithmic bias,''
\emph{IBM Journal of R\&D}, vol.~63, no.~4/5, 2019.

\bibitem{athey2019cate}
S.~Athey and G.~W.~Imbens, ``Machine learning methods for estimating
heterogeneous causal effects,'' \emph{Annu.\ Rev.\ of Econ.}, 2019.

\bibitem{econml2019}
Microsoft Research, ``EconML: A Python package for ML-based heterogeneous
treatment effects estimation,'' v.\ 0.x, 2019.

\bibitem{ibm2018kenexa}
IBM Corp., ``Kenexa Predictive Attrition,'' product brief, 2018.

\bibitem{workday2022retention}
Workday Inc., ``Workday Retention Risk,'' product brief, 2022.

\bibitem{microsoft2023viva}
Microsoft Corp., ``Microsoft Viva Insights,'' product brief, 2023.

\bibitem{visier2024}
Visier Inc., ``People analytics platform,'' product brief, 2024.

\bibitem{crunchr2024}
Crunchr B.V., ``HR analytics dashboards,'' product brief, 2024.

\bibitem{dwork2014algorithmic}
C.~Dwork and A.~Roth, ``The algorithmic foundations of differential
privacy,'' \emph{Foundations and Trends in TCS}, vol.~9, 2014.

\bibitem{lewis2020rag}
P.~Lewis et al., ``Retrieval-augmented generation for knowledge-intensive
NLP tasks,'' \emph{NeurIPS}, 2020.
\end{thebibliography}

\end{document}
